%!TEX root = ../main.tex
%%=====================================================================
%%  §1  INTRODUCTION          target 4-5 pp.
%%
%%  REWRITTEN 2026-08-21 on the v3 panel, to the JFEC house sequence
%%  (motivation -> literature -> gap/design -> contributions -> results
%%  -> roadmap). Two structural changes from the 2026-08-16 version:
%%   * the literature positioning is UNCOMMENTED and moved up. 9 of 9
%%     sampled JFEC introductions carry literature inside the intro.
%%   * contributions moved from 0.85 of the way in to 0.35, and turned
%%     from a nested enumerate into running ordinals. 0 of 9 sampled
%%     introductions use a list of any kind.
%%  NUMERALS: deliberately almost none. 5 of 9 sampled introductions
%%  report no numeric finding at all, and half of the figures the old
%%  version printed are either superseded on v3 or not yet independently
%%  recomputed. Every one of them now lives in the section or table this
%%  text points at. Do not reintroduce a numeral here without checking
%%  it against the exhibit that carries it.
%%  Voice: the author edits.
%%
%%  POSITIONING (unchanged, still binding):
%%   * brutti2026orthogonal ALREADY have a no-clustering control. Ours
%%     CORROBORATES; claim no priority. What their design cannot deliver
%%     is the moment decomposition (their P&L is normalized by its own
%%     standard deviation) and the partition-randomizing null (their
%%     placebo permutes within clusters, partition fixed).
%%   * guijarro2025deep report SR, mu and sigma per model. "We report
%%     volatility separately" is NOT a contribution.
%%
%%  FORBIDDEN HERE (each tested, each failed):
%%   - "every clustering earns the same alpha" (ANOVA rejects)
%%   - "the best method earns less alpha than no clustering"
%%   - any era rank-stability claim, in EITHER direction
%%   - any adjacent-pair ordering, any tradability claim
%%   - any point value for the head lift (bound only, +0.74 to +1.50)
%%   - any post-2009 alpha edge over K=1
%%   - "a clusterizer never sees a payoff" (false for four rules)
%%   - "forces/imposes group neutrality" (refuted, §4.3)
%%=====================================================================
\section{Introduction}
\label{sec:intro}

Machine learning methods are now used at most stages of trading
strategies \citep{kelly2023financial}, and clustering has become one of
the standard ways to form groups of comoving stocks for relative-value trading. First,
a clustering-based statistical arbitrage strategy partitions the stocks in the cross-section into several groups.
Then, a trading model builds trading portfolios on that partition, typically from returns hedged against the market. Previous studies in clustering-based trading usually evaluate the clustering algorithm and the trading model together and infer their effectiveness from the Sharpe ratio of the combined strategy \citep{Jinetal2023, LeeChang2026, brutti2026orthogonal}. However, the end-to-end performance does not identify which part makes what contribution. In this paper, we study what clustering contributes to statistical arbitrage. We call the clustering algorithm the clusterizer. First, we measure how much the clusterizer and the trading model contribute separately, and whether each contributes through the mean or through the volatilities of the payoffs. Second, we ask whether clusterizers based on different clustering rules recover different information. Third, we test whether traditional factor models, a reversal factor, or per-stock peer regressions, which regress each stock on the stocks most correlated with it, can provide the information that a clusterizer contributes.

The cluster structure of the stock market was studied long before cluster-based trading strategies were developed. The clusters can come from industry taxonomies, from the statistical structure of returns, and from a trading objective. Industry taxonomies fix the clusters in advance according to the economic characteristics of the companies
\citep{bhojraj2003whats}. The clusters can also be estimated from
returns, usually from the correlation matrix \citep{Mantegna1999}.
\citet{TOLA2008235} cluster the correlation matrix for portfolio construction.
\citet{Jinetal2023} apply the same idea to statistical arbitrage. Most recently, \citet{LeeChang2026} jointly fit
the partition and the model that trades based on it, via a transformer trained on a trading-based objective function.

A neighboring strand of literature asks how we can utilize the partitions for trading.
Early studies of relative-value trading include \citet{Gatevetal2006}, who documented the profitability of trading a single mean-reverting spread. Later, \citet{avellaneda2010statistical} replaced stock pairs with the factor-model residual, the part of a stock's return left after systematic factors are removed. It is worth noting that a cluster also provides another such factor, via the cluster mean. More recently, \citet{Jinetal2023}, \citet{LeeChang2026}, and
\citet{brutti2026orthogonal} each propose clusterizers of their own together with a strategy that gives trading portfolios based on their partitions. \citet{Jinetal2023} compare five spectral clusterizers, four of which are signed clusterizers that separate positive from negative correlations, under the same trading rule. In contrast, we separately analyze the contributions of clustering and of the trading rule within the whole trading pipeline.

To the best of our knowledge, no previous study compares clusterizers from different families under a common trading model, or separates what the partition contributes from what the model contributes. When the clusterizers and the trading rules are evaluated together, a higher Sharpe ratio may come from a better partition, from a better model, or from the way that the two match, so the end-to-end performance cannot help us examine the contribution of each component. Attribution also requires examining the mean and the volatility of the payoff separately. \citet{guijarro2025deep} examine both for factor-based statistical arbitrage. They report the Sharpe ratio, mean, and volatility of every model they fit and vary the factor construction under different signal models. However, they did not include clustering-based statistical arbitrage in their close examination.

We examine the clustering stage and the trading stage separately.
For the clustering stage, we compare 12 clusterizers from six families (geometric, hierarchical, eigenstructure-based, learned, trading-driven, and a sector taxonomy) against two controls.
The no-clustering control ($\K=1$) puts all stocks in a single group. The random-partition control is a uniform random partition with the same number of groups as the data-driven partitions. We call the partition of a clusterizer a real partition to contrast it with the random partition. Among the six families, we propose the trading-driven family and the learned family for the first time. In the trading-driven family, a clusterizer chooses its groups to optimize the Sharpe or Sortino ratio of the portfolio that a fixed trading rule builds on these groups. The learned family consists of \textsc{sssnet} \citep{he2022sssnet}, a signed-graph neural network that we bring to statistical arbitrage for the first time. We compare what the families' partitions contribute under a common downstream pipeline, without attempting to estimate the best attainable performance of each family of clusterizers.

For the trading stage, we trade every partition in a rule-based component and in a trained component. Each trading component builds a portfolio of long and short positions on a partition, and we call the daily return of this portfolio the payoff. The rule-based component uses a fixed cross-sectional rule adapted from \citet{Jinetal2023}. The trained component uses a neural trading head in the design of \citet{guijarro2025deep}, with the long convolution of \citet{fu2023simple} in place of their transformer. Because the fixed rule has no fitted parameter, the differences in performance among the rule-based portfolios come from the clustering stage alone. We then use the trained trading component to show how the conclusions change when a fitted model trades based on the partitions. We run both trading components on the same panel of daily returns for the 500 largest US equities over 27 years.

To attribute performance to the clustering stage and the trading stage, we compare the rule-based and trained components and use the two controls, replications under different random seeds, and a decomposition of the Sharpe gap. First, the two controls identify what the clustering stage contributes. The $\K=1$ control measures the cost of a missing partition. The random control measures the cost of wrong partitions. Second, the replications under different random seeds show whether different clusterizers can be distinguished in their partitions, positions, or payoffs. We compare partitions by the adjusted Rand index (\textsc{ari}), a measure of the agreement between two partitions corrected for chance. We judge gaps in the Sharpe-ratio ranking of the clusterizers against the dispersion of the design's own replications. In a head-seed rerun, the trading head is trained again from a different random initialization (a different head seed). In a partition draw, a stochastic clusterizer's daily partitions are generated again with a different random seed (a different partition seed). Third, the decomposition identifies whether a Sharpe gap arises through the mean payoff or through the payoff volatility. We split the gap
into the mean payoff part and the volatility part (\secref{sec:split}), because the same sample
estimates volatility more precisely than the mean (\secref{sec:resolution}).

We find that clustering lowers the volatility of the portfolio and raises its alpha, the intercept of a regression of the portfolio's payoff on the Fama--French five factors and momentum. First, demeaning market-hedged returns within the groups of a real partition removes the comovement that the stocks in each group share. With that comovement removed, the portfolio of a real partition is $39\%$ to $42\%$ less volatile than either control in the rule-based component and $33\%$ to $35\%$ less volatile in the trained component (\secsref{sec:hedge}{sec:mechanism}). Most of the Sharpe gap between a real partition and either control arises through this lower volatility (\secref{sec:split}). Second, demeaning within the groups of a real partition gives a more informative relative-value input than demeaning against the whole cross-section or within a random group. Consistent with this, we find that the two controls have the lowest alphas of the 14 rule-based portfolios (\secref{sec:hedge}).

With a fixed partition and at the same gross exposure, the trading head raises alpha in all 14 portfolios and raises volatility by proportionally less (\secref{sec:head}). Although the head widens the alpha advantage over both controls, especially over the random control, it narrows the volatility reduction. A wrong partition therefore has a larger disadvantage in the trained component than in the rule-based component, in the sense that the random control falls further behind the clusterizers in alpha.

In the trained component, the Sharpe-ratio ranking of the 12 clusterizers largely agrees with their ranking in the rule-based component. The ranking therefore reflects the quality of the partition as a trading input, whichever trading component trades it. However, some clusterizers in the middle of the ranking move by up to five places. The two controls also change order. In the rule-based component, the random control earns a higher Sharpe ratio than the $\K=1$ control. In the trained component, the two controls earn the same Sharpe ratio on average across head seeds (\secref{sec:head}).

We also find that the 12 clusterizers, which cluster stocks by qualitatively different principles, produce partitions that agree little by the \textsc{ari}. However, the portfolios built based on these partitions are highly correlated. The partitions are more alike in what demeaning within their groups removes than in how they group the stocks. For every clusterizer, demeaning within its groups removes much more comovement than the random control removes. By construction, the $\K=1$ control does not remove any comovement. Ten of the 12 clusterizers remove nearly the same share of the cross-sectional variance of market-hedged returns
(\secref{sec:mechanism}). The similarity between clusterizers also increases at each later stage of the pipeline, in the positions that the partitions induce and again in
the payoffs that those positions earn. In the payoffs, two different clusterizers differ
about as much as one clusterizer differs from its own head-seed rerun (\secref{sec:gradient}). Thus, the data clearly distinguish a real partition from the no-clustering and random controls, but provide much weaker separation among the clusterizers themselves. To the extent that the clusterizers' Sharpe ratios differ, they differ mainly through the volatility of the portfolios.

Most adjacent gaps in the Sharpe ranking are smaller than the variation of the gap itself across head-seed reruns (\secref{sec:field}). A stochastic clusterizer is unstable even across its own partition draws. Specifically, two partition draws of the same clusterizer disagree on group assignments nearly as much as two different clusterizers do. Moreover, the trading head enlarges the change in the Sharpe ratio that the redraw causes (\secref{sec:pseed}). Because the portfolios hold highly similar positions, re-selecting each year the clusterizer with the highest trailing Sharpe ratio does not outperform the equal-weight portfolio of all 12 clusterizers (\secref{sec:gradient}).

In contrast, changing the estimation window or the number of groups can change the Sharpe ratio by an amount of the same order as the gap between the best and the worst clusterizer. The estimation window is the trailing window of returns hedged against the market from which a clusterizer forms its groups. Between the best and the worst clusterizer, the gap in Sharpe ratio is $0.83$. We lengthen the estimation window from one month to one year. Over that range, the Sharpe ratio of unsigned spectral clustering rises by $0.88$, while that of average linkage is the same at both ends. We also reduce the number of groups from the $30$ used throughout to $10$. This change lowers the Sharpe ratio by up to $0.36$ (\secref{sec:window}).

Neither a factor model nor per-stock peer regressions make a partition redundant. \citet{guijarro2025deep} remark that latent factor estimation can, under appropriate assumptions, take the place of clustering. Taken as stated, the remark would make the clustering stage redundant, so we test it. We project the leading principal components, and separately the Fama--French factors, out of the day's cross-section of market-hedged returns. For each partition, we measure its variance share, the share of the cross-sectional variance of the day's market-hedged returns that demeaning within its groups removes. After either projection, every real partition loses at most $56\%$ of its variance share. No real partition falls to the variance share that a random partition would remove (\secref{sec:lens}).

We then test whether known return factors explain these alphas. We regress the payoff of every portfolio on the Fama--French five factors plus momentum and then add a short-term reversal factor and the five leading principal components of the market residuals. Every portfolio, including the two control portfolios, keeps a positive and significant alpha. The reversal factor explains at most $42\%$ of any portfolio's alpha, and it explains a larger share of alpha for the control portfolios than for the clusterizers' portfolios. To measure what a real partition adds, we also regress each clusterizer's portfolio and the random control's portfolio on the portfolio of the $\K=1$ control, which is the same pipeline run without clustering. Each clusterizer's portfolio keeps $41\%$ to $68\%$ of its alpha, whereas the random control's portfolio keeps $8\%$ to $9\%$ (\secref{sec:spanning}).

Because clustering-based trading is a special case of peer-based trading, in which every stock in a group shares the same peer set, we also compare it with per-stock peer regressions, in which each stock has its own peer set. In a per-stock peer regression, a stock's market-hedged return is regressed on those of the stocks most correlated with it, and the trading head trades the part that the regression does not explain. We build two such portfolios, which differ only in the number of days over which the regression is fitted. We call them the peer-regression portfolios. Peer regressions are a stronger substitute for a partition than the factors are.
In the cross-section, they absorb much of a partition's variance share, considerably more than the principal components or the Fama--French factors absorb. The peer regressions do not, however, fully absorb what the partitions earn. Against the peer-regression portfolios and the factors together, both controls lose their alpha. In contrast, most of the 12 clusterizers keep their alpha. Conversely, regressed on the equal-weight portfolio of all 12 clusterizers, both peer-regression portfolios keep their alpha (\secref{sec:peer}).

Finally, at a trading cost of five basis points per dollar traded, every daily-rebalanced portfolio, including the two control portfolios, has a negative net return (\secref{sec:costs}).

The rest of the article is organized as follows. In \secref{sec:lit}, we review the related literature. In \secref{sec:design}, we describe the panel, the 12 clusterizers, the two controls, and the rule-based and trained components. In \secref{sec:riskreturn}, we measure how clustering changes the volatility and the alpha of the portfolios and show how demeaning within groups lowers the volatility. We then split the Sharpe gap between the clusterizers and the controls into a mean part and a volatility part and measure what the trading head adds to the performance. In \secref{sec:invariance}, we compare the differences among the 12 clusterizers with the variation across head-seed reruns and partition draws. In \secref{sec:factor}, we test whether factor models or portfolios built on per-stock peer regressions explain the alphas of the clusterizers' portfolios. In that section, we also measure the effect of transaction costs. \secref{sec:conclusion} concludes. The appendix describes the construction of each clusterizer and the architecture and training schedule of the trading head. It also reports the additional results cited in the text.
